% Minimal retained preamble for lane-p2 cards taken from
% sources/sigma-2010-081-generalized-complex-reduction/pinned-original-source.tex
% (Singular Reduction of Generalized Complex Manifolds, DOI 10.3842/SIGMA.2010.081).
%
% Every declaration below is copied verbatim from the author's own file:
%   lines 9-19  : the author's theorem declarations
%   line  21    : \numberwithin{equation}{section}
%   line  23    : \usepackage[all]{xy}
%   lines 26-36 : the author's own macro definitions
% The published paper uses the publisher class \documentclass{sigma} (chosen in
% the \ifpdf block at lines 1-6), which the assembler does not accept -- only
% article and amsart are supported, see NEW_TASK_AUTHORING section 16.1.  The
% wrapper therefore supplies the article class and the base packages
% amsmath/amssymb/amsthm/graphicx.  Nothing was renamed, redefined, added or
% removed; no macro was invented.  xy is retained because the retained body
% contains the author's own commutative diagram.
\documentclass{article}
\usepackage{amsmath,amssymb,amsthm,graphicx}
\usepackage[all]{xy}

\newtheorem{Theorem}{Theorem}[section]
\newtheorem{Lemma}[Theorem]{Lemma}
%\newtheorem{problem}[Theorem]{Problem}
%\newtheorem{conv}[Theorem]{Convention}
\newtheorem{Proposition}[Theorem]{Proposition}
\newtheorem{Corollary}[Theorem]{Corollary}
{\theoremstyle{definition}
\newtheorem{Remark}[Theorem]{Remark}
\newtheorem{Definition}[Theorem]{Definition}
\newtheorem{Example}[Theorem]{Example}
}

\numberwithin{equation}{section}

\newcommand{\deff}[1]{\textbf{\emph{#1}}}
\newcommand{\func}[3]{#1 \colon #2 \to #3}
\newcommand{\hookfunc}[3]{#1 \colon #2 \hookrightarrow #3}
\newcommand{\eval}[2]{\left. #1 \right|_{#2}}
\newcommand{\pont}[1]{\mathbb{T} {#1}}
\newcommand{\cpont}[1]{\mathbb{T}_{\mathbb{C}} {#1}}
%\newcommand{\tang}{\pmb{\top}}
\newcommand{\tang}{T}
\newcommand{\iprod}[2]{\left\langle {#1} , {#2} \right\rangle}
\newcommand{\iiprod}[2]{\left\langle \! \left\langle {#1} , {#2} \right\rangle \! \right\rangle}
\newcommand{\litlet}{\renewcommand{\labelenumi}{(\alph{enumi})}}
\title{Every component of the fixed point set of a canonical action is a split submanifold}
\author{Timothy E. Goldberg}
\date{doi:10.3842/SIGMA.2010.081; CC BY 4.0}
\begin{document}
\maketitle
\noindent\textit{Editorial note.} The counted unit is the article's own Proposition prop9 at source lines 382-384 together with the proof environment that begins on the next line after the statement (387-434); the pairing is the ordinary adjacent one.\par\medskip
\noindent\textit{Editorial note.} The retained context at source lines 138-141 is the article's own definition of Dirac and complex Dirac structures with the equivalent description of an almost GC structure; the retained context at source lines 312-314 is its own definitions of a splitting bundle and a split submanifold; the retained context at source lines 336-359 is its own definition of the lift of a group action to the generalized tangent bundle and of a canonical action, including the author's own commutative diagram.\par\medskip
\noindent\textit{Editorial note.} Licence: the licence evidence is not a JSON observation file for this source. The card relies on the article-specific attribution notice retained with the corpus at editable-corpus/licenses/2819-source-license.txt, which names the article, the DOI, the publisher and CC BY 4.0, together with the retained CC BY 4.0 legal code at editable-corpus/licenses/CC-BY-4.0.txt; the paths are given relative to the corpus root. No acceptance record was written and no new observation was made.\par\medskip
\noindent\textit{Editorial note.} The published paper is typeset with the publisher class sigma, which the assembler does not accept (article and amsart only). The retained preamble supplied with this task therefore reproduces, verbatim, the author's theorem declarations at source lines 9-19, his numberwithin line at source line 21, his xy package line at source line 23 and his macro definitions at source lines 26-36, wrapped in the supported article class with amsmath, amssymb, amsthm and graphicx. Nothing was renamed, redefined or added. xy is needed because the retained definition of a canonical action contains the author's diagram.\par\medskip
\noindent\textit{Editorial note.} No retained line contains a citation, so the material carries no bibliography. The retained lines contain no cross-reference either: the sentence of the article that refers back to its Proposition prop4 for the description of an almost GC structure is deliberately not retained, and the equivalent author formulation at source line 140 is retained instead.\par\medskip
\noindent\textit{Editorial note.} The article's main theorem MainThm, its corollary MainThm2 and its other propositions, lemmas and examples are not counted here and are not used as premises.\par\medskip
\noindent\textit{Allowed standard prerequisites (editorial).}\begin{itemize}
\item Almost GC structure as a complex Dirac structure. A maximally isotropic linear subbundle of the generalized tangent bundle T M is a Dirac structure and a maximally isotropic linear subbundle of its complexification is a complex Dirac structure; an almost GC structure on M is a complex Dirac structure E with E intersected with its conjugate equal to zero, and the same object may be written as a bundle map J on the generalized tangent bundle; the article's own formulation is retained at source lines 138-141.. Fixes the object supplied by the hypothesis of the counted proposition, namely the almost GC structure J on M, and the bundle map J that the proof averages against.
\item Splitting bundle and split submanifold. For a submanifold S of an almost GC manifold (M,J), a splitting bundle is a subbundle N of the restricted generalized tangent bundle over S with T M restricted to S equal to T S plus N, such that T S plus the annihilator of N inside the generalized tangent bundle is invariant under J; a submanifold admitting such a bundle is a split submanifold of (M,J); retained at source lines 312-314.. It is exactly the conclusion of the counted proposition: the proof constructs such an N by averaging over the compact group and then verifies the required J-invariance.
\item Canonical action of a Lie group on an almost GC manifold. A smooth action of a Lie group G on a manifold M lifts to an action on the generalized tangent bundle by the bundle automorphisms (g\_*,(g\textasciicircum{}\{-1\})\textasciicircum{}*); the action on an H-twisted GC manifold is canonical when H is G-invariant and the action on the generalized tangent bundle commutes with J, that is the corresponding diagram commutes for every g in G; retained at source lines 336-359.. Supplies the second hypothesis of the counted proposition and the commutation J composed with the group action that the proof uses to keep the fixed-point data invariant.
\end{itemize}\medskip
\begin{Remark}
Let $M$ be a manifold.  A maximally isotropic linear subbundle of $\pont{M}$, respectively~$\cpont{M}$, is called a \deff{Dirac structure}, respectively \deff{complex Dirac structure} on $M$.
Thus an almost GC structure on $M$ is a complex Dirac structure $E \subset \cpont{M}$ satisfying $E \cap \overline{E} = \{0\}$.
\end{Remark}

\begin{Definition}
Let $(M, \mathcal{J})$ be an almost GC manifold, and let $S \subset M$ be a submanifold.  A \deff{splitting bundle} for $S$ with respect to $(M, \mathcal{J})$ is a subbundle $N$ of $\tang M|_S \to S$ such that $\tang M|_S = \tang S \oplus N$ and $\tang S \oplus \operatorname{Ann} (N) \subset \pont{M}$ is invariant under $\mathcal{J}$. If a splitting bundle exists for~$S$, then $S$ is called a \deff{split submanifold} of $(M,\mathcal{J})$.
\end{Definition}

\begin{Definition}
Let $M$ be a manifold, and let $G$ be a Lie group acting smoothly on $M$.  This lifts to an action of $G$ on $\pont{M}$ by bundle automorphisms, given by
\[
  \func{\left( g_*, (g^{-1})^* \right)}{\pont{M}}{\pont{M}}
\]
for each $g \in G$, where $g_{\ast}$ is the {\bf pushforward} of tangent vectors by the map $\func{g}{M}{M}$ and $(g^{-1})^{\ast}$ is the {\bf pullback} of tangent covectors by the map $\func{g^{-1}}{M}{M}$.

Let $\mathcal{J}$ be an $H$-twisted GC structure on $M$.  We say that the $G$-action on $(M, \mathcal{J}, H)$ is \deff{canonical} if the following hold.
\begin{enumerate}\itemsep=0pt
  \item The dif\/ferential form $H$ is $G$-invariant, i.e.\ $g^* H = H$ for all $g \in G$.
  \item The action of $G$ on $\pont{M}$ commutes with $\mathcal{J}$, i.e.\ the diagram
\[
  \xymatrix{
  \pont{M} \ar[d]_{ \left( g_*, (g^{-1})^* \right) } \ar[r]^{\mathcal{J}}
  &
  \pont{M} \ar[d]^{ \left( g_*, (g^{-1})^* \right) } \\
  \pont{M} \ar[r]_{\mathcal{J}}
  &
  \pont{M}
  }
\]
commutes for all $g \in G$.
\end{enumerate}
\end{Definition}

\begin{Proposition}  \label{prop9}
Let $M$ be a manifold, and let $\mathcal{J}$ be an almost GC structure on $M$. Suppose the compact Lie group $G$ acts canonically on $(M,\mathcal{J})$.  Then each component of $M^{G}$ is a split submanifold of $(M,\mathcal{J})$.
\end{Proposition}

\begin{proof}
Let $(M^G)'$ be a component of $M^G$.
First, recall that for each $x \in (M^G)'$ the derivative of the action of $G$ at $x$ def\/ines a linear action of $G$ on $\tang_x M$, and that $\tang_x (M^G)' = (\tang_x M)^G$.  Let $dg$ be a bi-invariant Haar measure on $G$, adjusted so that $dg(G) = 1$.  Def\/ine a bundle map $\func{\pi}{(\tang M)|_{(M^G)'}}{\tang (M^G)'}$ by setting
\[
  \pi_x(v) := \int_G (g \cdot v) \, dg
  \qquad
  \text{for all} \quad v \in V,
\]
for each $x \in (M^G)'$.  Def\/ine the subbundle $N \subset (\tang M)|_{(M^G)'}$ by setting $N_x := \ker \pi_x$ for each $x \in \left( M^G \right)'$.  That $\pi$ is a bundle map and $N$ is a vector bundle follow from the naturality of the technique of averaging by integration.  Note also that $\tang_x M = (\tang_x M)^G \oplus N_x = \tang_x (M^G)' \oplus N_x$ for each $x \in \left( M^G \right)'$.  To conclude that $N$ is a splitting for $(M^G)' \subset (M,\mathcal{J})$, it remains only to show that $\tang (M^G)' \oplus \operatorname{Ann}(N) = (\tang M)^G|_{(M^G)'} \oplus \operatorname{Ann}(N)$ is preserved by $\mathcal{J}$.

Fix $x \in \left( M^G \right)'$, let $V = \tang_x M$, and let $W = N_x$.  We claim that $\operatorname{Ann}(W) = (V^*)^G$, which would imply that $(V \oplus V^*)^G = V^G \oplus (V^*)^G = V^G \oplus \operatorname{Ann}(W)$.  Since $\mathcal{J}_x$ commutes with the action of $G$ on $V \oplus V^*$, we know that $\mathcal{J} \left( (V \oplus V^*)^G \right) = (V \oplus V^*)^G$, and hence a proof of this claim completes the proof of this proposition.

Let $\lambda \in \operatorname{Ann}(W)$, let $g \in G$, and let $u \in V$.  Decompose $u$ as $u=v + w$ for $v \in V^G$ and $w \in W$.
Then
\begin{alignat*}{2}
  (g \cdot \lambda)(u)
  &=
  \lambda ( g^{-1} \cdot u)
  && \\
  &=
  \lambda ( g^{-1} \cdot v) + \lambda(g^{-1} \cdot w)
  && \\
  &=
  \lambda (v) + \lambda(g^{-1} \cdot w)
  && \text{since } v \in V^G \\
  &=
  \lambda (v) + 0
  && \text{since $w \in W$, $W$ is $G$-stable, and }
  \lambda \in \operatorname{Ann}(W) \\
  &=
  \lambda (v) + \lambda (w)
  && \text{since $w \in W$ and } \lambda \in \operatorname{Ann}(W) \\
  &=
  \lambda (u).
  &&
\end{alignat*}
Therefore $\lambda \in (V^*)^G$, so $\operatorname{Ann}(W) \subset (V^*)^G$.  Note that
\[
  \operatorname{Ann}(W)
  \cong
  (V / W)^*
  =
  \left( (V^G \oplus W) / W \right)^*
  \cong
  (V^G)^*.
\]
Hence $\dim \operatorname{Ann}(W) = \dim (V^G)$, so $\operatorname{Ann} (W) = (V^*)^G$.
\end{proof}

\end{document}
